% GENERATED FILE. Edit canonical lessons, then run npm run generate:books.
% canonical-source-hash: fnv1a64:66dd4b0c60a93614
% canonical-lessons: GE-C101-dreissig, GE-C101-vierzig, GE-C101-fuenfzig, GE-C101-hundert, GE-C101-tausend

\chapter{Thirty, Forty, Fifty, Hundred, Thousand}
\label{ch:ge-thirty-forty-fifty-hundred-thousand}
\hlchaptermodality{101}

\begin{chapteropening}
\textbf{By the end of this chapter:} \emph{I can say thirty, forty, fifty, a hundred and a thousand.}

Complete the last lesson of chapter 101: I can say thirty, forty, fifty, a hundred and a thousand.
\end{chapteropening}

\section[dreißig]{dreißig — thirty}

\label{lesson:GE-C101-dreissig}



\begin{quote}
Before the new one: say the German for the world, then the German for a language.
\end{quote}

\subsection*{You'll want to know: dreißig}
\textbf{dreißig} — \textquotedblleft{}thirty\textquotedblright{}.

\begin{cousinweb}
\textbf{dreißig} $\leftarrow$ Old High German \emph{drīzzug} $\leftarrow$ Proto-Germanic \emph{*þrīz tigiwiz}, \textquotedblleft{}three tens\textquotedblright{} — the same build as English \emph{thirty}. (EtymDB 2.1)
\end{cousinweb}

\begin{grammarlens}[title={What you've built}]
One more word for this chapter.
\end{grammarlens}

\subsection*{Your turn}
\begin{itemize}
\raggedright
  \item \emph{Say it:} \emph{dreißig}
  \item \emph{Say it:} \emph{dreißig}, once more
  \item \emph{Say it:} say \emph{dreißig}
  \item \emph{Recall:} say the German for the world, then the German for a language, then say \emph{dreißig} again
\end{itemize}

\begin{tcolorbox}[breakable,colback=teal!4,colframe=teal!35!black,title={Before you move on}]
What does dreißig mean? (\textquotedblleft{}Thirty\textquotedblright{}.) Say it once more.
\end{tcolorbox}
\section[vierzig]{vierzig — forty}

\label{lesson:GE-C101-vierzig}



\begin{quote}
Before the new one: say the German for a language, then the German for thirty.
\end{quote}

\subsection*{You'll want to know: vierzig}
\textbf{vierzig} — \textquotedblleft{}forty\textquotedblright{}.

\begin{cousinweb}
\textbf{vierzig} $\leftarrow$ Old High German \emph{fiorzug} $\leftarrow$ Proto-Germanic \emph{*fedwōr tigiwiz}, \textquotedblleft{}four tens\textquotedblright{} — the same build as English \emph{forty}. (EtymDB 2.1)
\end{cousinweb}

\begin{grammarlens}[title={What you've built}]
One more word for this chapter.
\end{grammarlens}

\subsection*{Your turn}
\begin{itemize}
\raggedright
  \item \emph{Say it:} \emph{vierzig}
  \item \emph{Say it:} \emph{vierzig}, once more
  \item \emph{Say it:} say \emph{vierzig}
  \item \emph{Recall:} say the German for a language, then the German for thirty, then say \emph{vierzig} again
\end{itemize}

\begin{tcolorbox}[breakable,colback=teal!4,colframe=teal!35!black,title={Before you move on}]
What does vierzig mean? (\textquotedblleft{}Forty\textquotedblright{}.) Say it once more.
\end{tcolorbox}
\section[fünfzig]{fünfzig — fifty}

\label{lesson:GE-C101-fuenfzig}



\begin{quote}
Before the new one: say the German for thirty, then the German for forty.
\end{quote}

\subsection*{You'll want to know: fünfzig}
\textbf{fünfzig} — \textquotedblleft{}fifty\textquotedblright{}.

\begin{cousinweb}
\textbf{fünfzig} is \emph{fünf} + \emph{-zig}, \textquotedblleft{}tens\textquotedblright{} — the same build as English \emph{fifty}. EtymDB 2.1 traces \emph{fünf} to Old High German \emph{fimf} $\leftarrow$ Proto-Germanic \emph{*fimf}.
\end{cousinweb}

\begin{grammarlens}[title={What you've built}]
One more word for this chapter.
\end{grammarlens}

\subsection*{Your turn}
\begin{itemize}
\raggedright
  \item \emph{Say it:} \emph{fünfzig}
  \item \emph{Say it:} \emph{fünfzig}, once more
  \item \emph{Say it:} say \emph{fünfzig}
  \item \emph{Recall:} say the German for thirty, then the German for forty, then say \emph{fünfzig} again
\end{itemize}

\begin{tcolorbox}[breakable,colback=teal!4,colframe=teal!35!black,title={Before you move on}]
What does fünfzig mean? (\textquotedblleft{}Fifty\textquotedblright{}.) Say it once more.
\end{tcolorbox}
\section[hundert]{hundert — a hundred}

\label{lesson:GE-C101-hundert}



\begin{quote}
Before the new one: say the German for forty, then the German for fifty.
\end{quote}

\subsection*{You'll want to know: hundert}
\textbf{hundert} — \textquotedblleft{}a hundred\textquotedblright{}.

\begin{cousinweb}
\textbf{hundert} $\leftarrow$ Old High German \emph{hundert} $\leftarrow$ Proto-Germanic \emph{*hundaradą}, built from \emph{*hundą}, \textquotedblleft{}hundred\textquotedblright{}, and \emph{*radą}, \textquotedblleft{}count\textquotedblright{} (EtymDB 2.1). English \emph{hundred} is the same word.
\end{cousinweb}

\begin{grammarlens}[title={What you've built}]
One more word for this chapter.
\end{grammarlens}

\subsection*{Your turn}
\begin{itemize}
\raggedright
  \item \emph{Say it:} \emph{hundert}
  \item \emph{Say it:} \emph{hundert}, once more
  \item \emph{Say it:} say \emph{hundert}
  \item \emph{Recall:} say the German for forty, then the German for fifty, then say \emph{hundert} again
\end{itemize}

\begin{tcolorbox}[breakable,colback=teal!4,colframe=teal!35!black,title={Before you move on}]
What does hundert mean? (\textquotedblleft{}A hundred\textquotedblright{}.) Say it once more.
\end{tcolorbox}
\section[tausend]{tausend — a thousand}

\label{lesson:GE-C101-tausend}



\begin{quote}
Before the new one: say the German for fifty, then the German for a hundred.
\end{quote}

\subsection*{You'll want to know: tausend}
\textbf{tausend} — \textquotedblleft{}a thousand\textquotedblright{}.

\begin{cousinweb}
\textbf{tausend} $\leftarrow$ Old High German \emph{dūsunt} $\leftarrow$ Proto-Germanic \emph{*þūsundī}, \textquotedblleft{}thousand\textquotedblright{} — inherited, the same word as English \emph{thousand}. (EtymDB 2.1)
\end{cousinweb}

\begin{grammarlens}[title={What you've built}]
One more word for this chapter.
\end{grammarlens}

\subsection*{Your turn}
\begin{itemize}
\raggedright
  \item \emph{Say it:} \emph{tausend}
  \item \emph{Say it:} \emph{tausend}, once more
  \item \emph{Say it:} say \emph{tausend}
  \item \emph{Recall:} say the German for fifty, then the German for a hundred, then say \emph{tausend} again
\end{itemize}

\begin{tcolorbox}[breakable,colback=teal!4,colframe=teal!35!black,title={Before you move on}]
What does tausend mean? (\textquotedblleft{}A thousand\textquotedblright{}.) Say it once more.
\end{tcolorbox}
